\documentclass[a4paper,12pt]{article}
\usepackage[utf8]{inputenc}
\usepackage[T1]{fontenc}
\usepackage[french]{babel}
\usepackage{geometry}
\geometry{left=2.5cm,right=2.5cm,top=2.5cm,bottom=2.5cm}

\usepackage{lmodern}
\usepackage{xcolor}
\definecolor{titleblue}{RGB}{0,51,140}
\definecolor{TITLEBLUE}{RGB}{0,51,140}

\usepackage{hyperref}

\pagestyle{empty}

\begin{document}

\begin{flushleft}
    \textbf{Ismail AISSAOUI} \\
    Ingénieur d'État en Intelligence Artificielle \\
    Chlef, Algérie \\
    ismail.aissaoui.pro@gmail.com \\
    +213 660 70 77 96 \\
    \href{https://ismail-aissaoui.vercel.app}{ismail-aissaoui.vercel.app}
\end{flushleft}

\begin{flushright}
    À l'attention du \\
    \textbf{Professeur Vincent CHEUTET} \\
    Laboratoire DISP \\
    INSA Lyon \\
    69621 Villeurbanne, France \\
    \medskip
    Le 29 Avril 2026
\end{flushright}

\bigskip

\noindent \textbf{Objet : Candidature pour un doctorat en Jumeaux Numériques \& Résilience des Systèmes}

\bigskip

Monsieur le Professeur,

C'est avec une grande admiration pour les travaux du laboratoire DISP, et tout particulièrement vos recherches sur l'articulation entre simulation et Jumeaux Numériques au service de la résilience industrielle (projet JUNEAU), que je vous adresse ma candidature pour un doctorat sous votre direction.

Actuellement en fin de cycle d'ingénieur d'État à l'ESI-SBA (Algérie), je réalise mon projet de fin d'études en collaboration avec Sonatrach sur le développement d'un \textbf{Jumeau Numérique (Architecture Générative Mamba-SSM)} pour la maintenance prédictive de turbines à gaz. Dans ce cadre, j'ai conçu une architecture duale combinant un modèle "Edge" (\textbf{C++ JIT xLSTM}) et une approche de substitut génératif pour optimiser la précision du diagnostic sous contrainte matérielle.

Mon parcours m'a permis de développer une expertise solide en modélisation hybride ainsi qu'en détection d'anomalies via des approches innovantes comme les \textbf{PINN} (Physics-Informed Neural Networks). Je suis convaincu que mon expérience pratique sur des infrastructures critiques et ma passion pour l'architecture des systèmes cyber-physiques s'alignent parfaitement avec la vision du DISP sur l'industrie du futur.

Je souhaiterais approfondir la question de la résilience dynamique des systèmes de systèmes (SoS), en explorant comment les Jumeaux Numériques peuvent non seulement prédire les chocs, mais aussi guider l'adaptation structurelle en temps réel.

Vous trouverez ci-joint mon CV détaillé. Je reste à votre entière disposition pour un entretien (en personne ou par visioconférence) afin d'échanger sur la possibilité d'intégrer votre équipe pour la rentrée 2026.

Je vous prie d'agréer, Monsieur le Professeur, l'expression de mes salutations les plus distinguées.

\bigskip

\begin{flushright}
    \textbf{Ismail AISSAOUI}
\end{flushright}

\end{document}
